\documentclass[11pt,a4paper]{article}
\newcommand{\labonaamwaarde}{labo 5}
\usepackage{syllabus}

\begin{document}

\syllabustitel{labo 5 -- \emph{Datavisualisatie I}}{Vertel het verhaal met plotly}

% ---------------------------------------------------------------------------
% Gedeeld door de syllabi van labo 3, 4 en 5. Pas dit hier aan, niet in de labobestanden.
\section{Overzicht}

Tijdens labo 3, 4 en 5 ben jij de nieuwe data-analist van een
(fictieve) jeugdclub voor voetbal en atletiek. De hoofdtrainer bezorgt je
drie exportbestanden uit de systemen van de club. Die zijn bewust
rommelig, zoals echte exports: inconsistente labels, sprinttijden in
milliseconden, ontbrekende waarden, een atleet die twee keer voorkomt en
een kolom met een komma als decimaalteken.

Elke sessie bouwt verder op de vorige:

\begin{description}[leftmargin=1.6cm, style=nextline]
  \item[Labo 3] \textbf{Van bestand tot nette tabel.} Inladen, verkennen,
    filteren, kolommen toevoegen en opkuisen tot je de data kunt
    vertrouwen.
  \item[Labo 4] \textbf{Combineren, samenvatten, herschikken.} Groepen
    samenvatten, tabellen koppelen, wisselen tussen een brede en een lange
    vorm, en werken met een hartslagreeks per seconde.
  \item[Labo 5] \textbf{Vertel het verhaal met plotly.} De juiste grafiek
    kiezen voor een vraag, eerlijke grafieken maken met
    \code{plotly.express} en ze combineren tot een klein interactief
    dashboard.
\end{description}


% ---------------------------------------------------------------------------
\section{Leerdoelen}

Na labo 5 kun je:

\begin{enumerate}[label=\textbf{LD\arabic*}, leftmargin=1.4cm]
  \item een grafiektype kiezen dat bij de vraag past, en misleidende
        grafieken herkennen, zoals een afgeknotte as;
  \item met \code{plotly.express} histogrammen, boxplots,
        spreidingsdiagrammen en lijngrafieken maken, met kleur, facetten
        en grootte voor extra dimensies;
  \item grafieken interactief maken met UI-elementen van marimo en ze
        combineren tot een dashboard.
\end{enumerate}

% ---------------------------------------------------------------------------
\section{Voorkennis}

Labo 3 en 4 (\emph{Dataverwerking I en II}): data inladen, filteren en
opkuisen met \code{pandas}, groeperen met \code{groupby}, tabellen
koppelen met \code{merge} en het verschil tussen een lange en een brede
vorm. Labo 5 start met de volledige opkuis van labo 3 en 4 in één cel.

% ---------------------------------------------------------------------------
% Gedeeld door de syllabi van labo 3, 4 en 5. Pas dit hier aan, niet in de labobestanden.
\section{Software en lesmateriaal}

Elk labo is een \textbf{marimo-notebook} in de map \code{labs/} van de
cursusrepository. Je kunt het op twee manieren openen:

\begin{itemize}
    \item \textbf{Online:} surf naar \url{https://marimo.app}, kies
        \emph{New $\rightarrow$ Open from URL\dots} en plak de raw-URL van
        het notebook op GitHub. De data wordt dan van GitHub gelezen, dus
        je hoeft niets te installeren.
    \item \textbf{Lokaal:} installeer de pakketten uit
        \code{requirements.txt} en open het notebook met marimo. De data
        wordt dan uit de map \code{data/} gelezen.

\end{itemize}

Drie dingen die je moet weten over marimo:

\begin{itemize}
  \item Notebooks zijn \textbf{reactief}: als je een cel aanpast, worden
        alle cellen die ervan afhangen opnieuw uitgevoerd, in de volgorde
        van hun afhankelijkheden.
  \item \textbf{Twee cellen kunnen niet dezelfde variabele definiëren.}
        Geef elke stap een eigen, duidelijke naam (\code{tests\_raw},
        \code{tests\_clean}, \dots) of zet een reeks stappen samen in één
        cel.
  \item Elke oefening bevat een placeholder \code{...}. Een verborgen cel
        eronder controleert je antwoord: die kleurt \textbf{groen} als het
        juist is, en geeft anders een algemene hint over wat nog niet
        klopt, zonder het antwoord te verklappen.
\end{itemize}


% ---------------------------------------------------------------------------
% Gedeeld door de syllabi van labo 3, 4 en 5. Pas dit hier aan, niet in de labobestanden.
\section{De dataset}

\noindent
\rowcolors{2}{tint}{white}
\begin{tabularx}{\textwidth}{@{}L{3.1cm}L{3.2cm}Y@{}}
\toprule
\textbf{Bestand} & \textbf{Omvang} & \textbf{Inhoud en gekende problemen} \\
\midrule
\code{athletes.csv}   & 43 rijen, 7 kolommen   & ID, naam, geboortedatum, geslacht, positie, lengte, gewicht. 13 schrijfwijzen voor 4 posities, gewicht als tekst met een decimale komma (\code{"57,5"}), één atleet twee keer ingevoerd met twee ID's, enkele ontbrekende waarden. \\
\code{tests.csv}      & 114 rijen, 8 kolommen  & Sprint over 10\,m en 30\,m, spronghoogte (countermovement jump) en Yo-Yo-niveau, op 3 testmomenten. Sommige sprinttijden staan in milliseconden; enkele ontbrekende waarden. \\
\code{hr\_session.csv} & 12\,000 rijen, 5 kolommen & Hartslag en snelheid per seconde van 5 atleten tijdens één training van 40 minuten. \\
\bottomrule
\end{tabularx}
\rowcolors{1}{}{}


% ---------------------------------------------------------------------------
\needspace{18\baselineskip}
\section{Planning}

% Planning van één labo. Gebruikt door de syllabus van dit labo en door de syllabus voor labo 3-5.
\textbf{Doel:} een grafiektype kiezen voor een gegeven vraag, de
standaardgrafieken van \code{plotly.express} maken, extra dimensies
toevoegen met kleur, facetten en grootte, en eenvoudige interactiviteit
inbouwen.

\begin{schedule}
13:00--13:30 & Principes & De juiste grafiek voor de vraag (verdeling, vergelijking, verband, evolutie); kritiek op een misleidende grafiek. \\
13:30--14:30 & Basis van \code{plotly.express} & Histogram, boxplot, spreidingsdiagram en lijngrafiek, met titels, aslabels en hover-informatie. \\
14:30--15:15 & Meer dimensies & \code{color}, \code{facet\_col} en \code{size}; één vaste kleur per positie. \\
15:15--15:30 & Pauze & \\
15:30--16:15 & Hartslag als tijdreeks & Zonebanden met \code{add\_hrect}, en waarom we geen tweede y-as gebruiken. \\
16:15--16:45 & Interactiviteit & \code{mo.ui.dropdown}, \code{mo.ui.range\_slider}, optioneel selecties met \code{mo.ui.plotly}. \\
16:45--17:30 & Minidashboard & 3 grafieken en 1 keuzemenu waarop minstens één grafiek reageert. \\
\end{schedule}

\textbf{Oefeningen:} (1) een histogram van het gewicht; (2) Yo-Yo-niveau
per positie, met facetten per geslacht; (3) snelheid doorheen de tijd
voor één atleet; \textbf{opdracht:} het minidashboard.


% ---------------------------------------------------------------------------
\section{Werkwijze}

\begin{itemize}
  \item Elk blok start met een korte demo in het notebook, gevolgd door
        oefeningen die je in het notebook zelf oplost.
  \item Gebruik de automatische controles als feedback, niet als doel:
        als een controle groen kleurt, zorg dan dat je kunt uitleggen
        \emph{waarom} je code werkt.
  \item Het minidashboard bouw je met alles uit dit labo: het doel is
        combineren, niet nieuwe grafiektypes leren.
\end{itemize}

% ---------------------------------------------------------------------------
\section{Evaluatie}

\begin{itemize}
  \item \textbf{Minidashboard:} drie grafieken en één interactief
        keuzemenu, gemaakt met de data van de club. Het kan dienen als je
        inzending voor deze module en bereidt de casepresentaties en het
        Project Lab van volgend semester voor.
\end{itemize}

% ---------------------------------------------------------------------------
\section{Voorbereiding}

\begin{itemize}
  \item Controleer vóór de sessie dat je het notebook kunt openen,
        lokaal of op marimo.app.
  \item Herhaal \code{groupby}, \code{merge} en de lange vorm uit labo 4.
        Grafieken met plotly express vertrekken van data in de lange vorm.
  \item Lees de cursustekst bij labo 5 naast het notebook.
\end{itemize}

\end{document}
